\documentclass[12pt,a4paper]{ctexart}
\usepackage{geometry}
\geometry{top=2.5cm,bottom=2.5cm,left=2.5cm,right=2.5cm}
\usepackage{amsmath,amssymb,mathtools}
\usepackage{enumitem,array,booktabs,pifont}
\usepackage{hyperref,fancyhdr}
\pagestyle{fancy}
\fancyhf{}
\fancyhead[C]{\textbf{数学 · 限时模拟卷三}}
\fancyfoot[C]{\thepage}
\title{\textbf{数学限时模拟卷三}\\[5pt]\large 满分150分\quad 时间120分钟}
\date{}

\begin{document}
\maketitle

\section*{各档次答题策略}
\begin{center}
\begin{tabular}{p{2.5cm}|p{4cm}|p{4cm}|p{4cm}}
\textbf{目标} & \textbf{必做题} & \textbf{尽力题} & \textbf{建议放弃} \\ \hline
\textbf{90分} & 1-6、9-10、13-14、17-18 & 7、19(1)、20(1) & 8、11-12、15-16、21-22 \\ \hline
\textbf{120分} & 1-7、9-11、13-15、17-20 & 8、16、21 & 12、22(2) \\ \hline
\textbf{140+分} & 全做 & 全做 & 无 \\
\end{tabular}
\end{center}

\section*{一、单选题（8×5=40分）}

\textbf{1.} $A=\{x\mid x^2-2x-3\le0\}$，$B=\{0,1,2,3,4\}$，$A\cap B=$
A. $\{0,1,2\}$ \quad B. $\{0,1,2,3\}$ \quad C. $\{1,2,3\}$ \quad D. $\{0,1,2,3,4\}$

\textbf{2.} $|z|$，$z=\frac{3+i}{1-2i}$，则$|z|=$
A. $2$ \quad B. $\sqrt2$ \quad C. $\sqrt5$ \quad D. $5$

\textbf{3.} $\sin\alpha=\frac35$，$\alpha\in(\frac{\pi}{2},\pi)$，$\cos(\alpha+\frac{\pi}{4})=$
A. $-\frac{7\sqrt2}{10}$ \quad B. $\frac{7\sqrt2}{10}$ \quad C. $-\frac{\sqrt2}{10}$ \quad D. $\frac{\sqrt2}{10}$

\textbf{4.} $f(x)=x^2-2\ln x$的单调递减区间为
A. $(0,1)$ \quad B. $(-\infty,-1)\cup(0,1)$ \quad C. $(-\infty,-1)$ \quad D. $(0,+\infty)$

\textbf{5.} $\vec{a}=(1,2)$，$\vec{b}=(2,k)$，夹角为锐角，则$k$范围是
A. $k>-1$ \quad B. $k>-1$且$k\neq4$ \quad C. $k<-1$ \quad D. $k\in\mathbb{R}$

\textbf{6.} 等差$\{a_n\}$中$a_1+a_4+a_7=15$，则$a_4=$
A. $3$ \quad B. $5$ \quad C. $7$ \quad D. $9$

\textbf{7.} 过$A(2,1)$与圆$x^2+y^2=5$相切的直线方程为
A. $2x+y=5$ \quad B. $x+2y=5$ \quad C. $2x+y-5=0$ \quad D. $x+2y-4=0$

\textbf{8.} $f(x)=a^x+\log_ax$在$[1,2]$上最大值与最小值之和为$\log_a2+6$，则$a=$
A. $\frac12$ \quad B. $2$ \quad C. $3$ \quad D. $4$

\section*{二、多选题（4×5=20分）}

\textbf{9.} 真命题是
A. $\exists x\in\mathbb{R}$，$x^2+2x+2<0$ \quad B. $\forall x>0$，$x+\frac1x\ge2$
C. 若$a>b$则$a^2>b^2$ \quad D. 若$a>b>0$则$\frac1a<\frac1b$

\textbf{10.} $f(x)=A\sin(\omega x+\varphi)$，$A>0,\omega>0,|\varphi|<\frac{\pi}{2}$，其图象过点$(0,1)$且$(\frac{\pi}{6},2)$为最高点，则
A. $A=2$ \quad B. $\omega=2$ \quad C. $\varphi=\frac{\pi}{6}$ \quad D. $f(\frac{\pi}{12})=\sqrt3$

\textbf{11.} 已知$m,n$为直线，$\alpha,\beta$为平面，则
A. $m\perp\alpha,n\subset\beta,m\perp n\Rightarrow\alpha\perp\beta$
B. $m\perp\alpha,n\perp\beta,\alpha\perp\beta\Rightarrow m\perp n$
C. $m\parallel n,n\subset\beta\Rightarrow m\parallel\beta$
D. $m\perp\alpha,n\subset\alpha\Rightarrow m\perp n$

\textbf{12.} $f(x)=x^3-3x+1$，则
A. 有2个极值点 \quad B. 极大值为3 \quad C. 图象关于$(0,1)$对称 \quad D. $f(x)=0$有3实根

\section*{三、填空题（4×5=20分）}

\textbf{13.} $\int_0^1(3x^2+2x)dx=$ \underline{\hspace{2cm}}。

\textbf{14.} 抛物线$y^2=4x$的焦点到准线距离为\underline{\hspace{2cm}}。

\textbf{15.} $a=0.3^{0.2}$，$b=0.2^{0.3}$，$c=\log_{0.3}0.2$，大小关系为\underline{\hspace{2cm}}。

\textbf{16.} $a_1=1$，$a_{n+1}=a_n+2n+1$，则$a_n=$ \underline{\hspace{2cm}}。

\newpage
\section*{四、解答题（共70分）}

\textbf{17.}（10分）等差$\{a_n\}$中$a_2=3$，$a_5=9$。（1）求通项；（2）$b_n=\frac1{a_na_{n+1}}$，求$T_n$。

\vspace{6cm}

\textbf{18.}（12分）$\triangle ABC$中$\sin^2A-\sin^2B-\sin^2C=\sin B\sin C$。（1）求$A$；（2）$a=2\sqrt3$，求$b+c$最大值。

\vspace{7cm}

\textbf{19.}（12分）三棱锥$P-ABC$中$PA\perp$面$ABC$，$\angle ABC=90^\circ$，$PA=AB=BC=2$。（1）证$BC\perp$面$PAB$；（2）求$PC$与面$PAB$所成角正弦值。

\vspace{7cm}

\textbf{20.}（12分）袋中4白3黑，每次取1个不放回，直到取到白球。$X$为取球次数。（1）求$X$分布列；（2）求$E(X)$。

\vspace{6cm}

\textbf{21.}（12分）椭圆$\frac{x^2}{4}+\frac{y^2}{3}=1$，$A(-2,0),B(2,0)$。过$P(1,0)$直线交椭圆于$M,N$，$AM$与$BN$交于$Q$。求$Q$轨迹方程。

\vspace{7cm}

\textbf{22.}（12分）$f(x)=\ln x-\frac12ax^2$。（1）讨论单调性；（2）若$f(x)$在$(0,+\infty)$上恒负，求$a$范围。

\newpage
\section*{参考答案与解析}

\subsection*{一、单选题}
\textbf{1. B} $A=[-1,3]$，$A\cap B=\{0,1,2,3\}$\\
\textbf{2. B} $z=\frac{3+i}{1-2i}=\frac{(3+i)(1+2i)}{1+4}=\frac{1+7i}{5}$，$|z|=\frac{\sqrt{1^2+7^2}}{5}=\frac{5\sqrt2}{5}=\sqrt2$\\
\textbf{3. C} $\cos\alpha=-\frac45$，$\cos(\alpha+\frac{\pi}{4})=\frac{\sqrt2}{2}(\cos\alpha-\sin\alpha)=\frac{\sqrt2}{2}(-\frac45-\frac35)=-\frac{7\sqrt2}{10}$\\
\textbf{4. A} 定义域$(0,+\infty)$，$f'(x)=2x-\frac2x=\frac{2(x^2-1)}{x}<0\Rightarrow0<x<1$\\
\textbf{5. B} $\vec{a}\cdot\vec{b}=2+2k>0\Rightarrow k>-1$，且不共线$\Rightarrow\frac12\neq\frac2k\Rightarrow k\neq4$\\
\textbf{6. B} $a_1+a_4+a_7=3a_4=15\Rightarrow a_4=5$\\
\textbf{7. A} 点$A$在圆上，切线$2x+y=5$\\
\textbf{8. B} $a>1$时$f$递增，$f(1)+f(2)=a+\log_a1+a^2+\log_a2=a+a^2+\log_a2=\log_a2+6\Rightarrow a^2+a-6=0\Rightarrow a=2$

\subsection*{二、多选题}
\textbf{9. BD} A中$\Delta=4-8<0$恒正；C反例$a=1,b=-2$\\
\textbf{10. ABC} $A=2$，$\frac{T}{4}=\frac{\pi}{6}\Rightarrow T=\frac{2\pi}{3}\Rightarrow\omega=3$。代入$(0,1)$得$\sin\varphi=\frac12\Rightarrow\varphi=\frac{\pi}{6}$\\
\textbf{11. BD} A反例$m\perp\alpha,n\parallel\alpha$；C需要$m\not\subset\beta$\\
\textbf{12. ACD} $f'(x)=3x^2-3=3(x-1)(x+1)$，$f(-1)=3$，$f(1)=-1$，极大3极小-1，3实根

\subsection*{三、填空题}
\textbf{13.} $[x^3+x^2]_0^1=1+1=2$\\
\textbf{14.} $p=2$，焦点到准线距离$p=2$\\
\textbf{15.} $c>0$，$0<a,b<1$，$c>1$，$c>a>b$（或写$c>a>b$）\\
\textbf{16.} $a_n=1+\sum_{k=1}^{n-1}(2k+1)=1+(n-1)+(n-1)^2=n^2$\\
（验证：$a_1=1$，$a_2=1+3=4=2^2$，$a_3=4+5=9=3^2$）

\subsection*{四、解答题}
\textbf{17.}（1）$d=\frac{a_5-a_2}{3}=2$，$a_n=a_2+(n-2)d=3+2(n-2)=2n-1$\\
（2）$b_n=\frac1{(2n-1)(2n+1)}=\frac12(\frac1{2n-1}-\frac1{2n+1})$，$T_n=\frac12(1-\frac1{2n+1})=\frac{n}{2n+1}$

\textbf{18.}（1）正弦定理：$a^2=b^2+c^2+bc$，$\cos A=\frac{b^2+c^2-a^2}{2bc}=-\frac12$，$A=120^\circ$\\
（2）$2R=\frac{a}{\sin A}=\frac{2\sqrt3}{\frac{\sqrt3}{2}}=4$，$b+c=2R(\sin B+\sin C)$，$C=60^\circ-B$\\
$=4[\sin B+\sin(60^\circ-B)]=4[\sin B+\frac{\sqrt3}{2}\cos B-\frac12\sin B]$
$=4(\frac12\sin B+\frac{\sqrt3}{2}\cos B)=4\sin(B+60^\circ)\le4$\\
当$B=30^\circ$时取最大值$4$

\textbf{19.}（1）$BC\perp AB$（已知），$PA\perp$面$ABC\Rightarrow PA\perp BC$，故$BC\perp$面$PAB$\\
（2）由（1）知$\angle CPB$即为$PC$与面$PAB$所成角。$PB=2\sqrt2$，$BC=2$，$PC=\sqrt{PB^2+BC^2}=2\sqrt3$，$\sin\angle CPB=\frac{BC}{PC}=\frac{2}{2\sqrt3}=\frac{\sqrt3}{3}$

\textbf{20.}（1）$P(X=1)=\frac47$，$P(X=2)=\frac37\cdot\frac45=\frac{12}{35}$，$P(X=3)=\frac37\cdot\frac15\cdot\frac44=\frac3{35}$\\
（2）$E(X)=1\cdot\frac47+2\cdot\frac{12}{35}+3\cdot\frac3{35}=\frac{20+24+9}{35}=\frac{53}{35}\approx1.514$

\textbf{21.} 设$MN:x=my+1$，$M(x_1,y_1)$，$N(x_2,y_2)$，联立得$(3m^2+4)y^2+6my-9=0$。\\
韦达定理得$y_1+y_2=-\frac{6m}{3m^2+4}$，$y_1y_2=-\frac{9}{3m^2+4}$。\\
$AM:y=\frac{y_1}{x_1+2}(x+2)$，$BN:y=\frac{y_2}{x_2-2}(x-2)$。联立消$y$代入$x_1=my_1+1$，$x_2=my_2+1$，利用韦达定理化简得$x=4$。\\
$Q$的轨迹方程为$x=4$（直线）

\textbf{22.}（1）$f'(x)=\frac1x-ax=\frac{1-ax^2}{x}$，$x>0$\\
$a\le0$时$f'(x)>0$，$(0,+\infty)$递增\\
$a>0$时$x\in(0,\frac1{\sqrt a})$增，$(\frac1{\sqrt a},+\infty)$减\\
（2）$f(x)<0$恒成立$\Rightarrow f_{\max}<0$。$a\le0$时$f(x)$无上界，不恒负。\\
$a>0$时$f_{\max}=f(\frac1{\sqrt a})=\ln\frac1{\sqrt a}-\frac12=-\frac12\ln a-\frac12=-\frac12(\ln a+1)<0\Rightarrow\ln a>-1\Rightarrow a>\frac1e$

\end{document}
